\section{Introduction}
\label{sec:tpc58-introduction}

The fixed-shift residual program has reached a point at which two
nonnegative quantities that look similar must be separated sharply.
The very-high-prime decomposition of TPC--45 produces the physical
all-minus affine corner
\begin{equation}
 \cN_{\rm nat}
 :=\left\|v_0-\sum_{p>y}v_p\right\|^2,
 \label{eq:tpc58-intro-native}
\end{equation}
whereas averaging the low Walsh signs produces
\begin{equation}
 \cR_{\rm low}
 :=\cD_0-\cD_H
   +\left\|\sum_{p>y}x_p\right\|^2.
 \label{eq:tpc58-intro-residual}
\end{equation}
For each individual \(p\), the diagonal norms agree,
\(\|v_p\|=\|x_p\|\), but the cross inner products do not
\citep{WangTPC45}.  The prime-resolved work in TPC--55--57 addresses the
coherent sum over \(p\) at a fixed residual label.  It does not yet
return that residual energy to the native output space
\citep{WangTPC55,WangTPC56,WangTPC57}.

The missing operation is simple to state and consequential to analyze:
\begin{equation}
 p\longmapsto(i,s)
 \quad\hbox{followed by}\quad
 (i,s)\longmapsto i.
 \label{eq:tpc58-two-groupings}
\end{equation}
The first arrow groups terminal primes while retaining the residual
label \(s\).  The second erases \(s\) inside the native output coordinate
\(i\).  TPC--57 studies the first arrow.  This paper studies the second.
In particular, one prime in every \((i,s)\)-fiber does not imply one
residual label in every \(i\)-fiber.

\subsection{The exact erasure identity}

Fix the baseline active residual labels \(\cL_i\) at output \(i\), and
write \(z_i=(z_{i,t})_{t\in\cL_i}\).  Once the M\"obius gauge is included,
the native map is
\begin{equation}
 (\Eop_\mu z)_i:=\sum_{t\in\cL_i}\mu(t)z_{i,t}.
 \label{eq:tpc58-intro-erasure}
\end{equation}
The entire discrepancy between resolved and native energies is
\begin{equation}
 \|\Eop_\mu z\|^2-\|z\|^2
 =2\operatorname{Re}\sum_i
   \sum_{t<t'}\mu(t)\overline{\mu(t')}
      z_{i,t}\overline{z_{i,t'}}.
 \label{eq:tpc58-intro-defect}
\end{equation}
For \(k_i=|\cL_i|>0\), the \(i\)-block has one nonzero singular value
\(\sqrt{k_i}\) and a \(k_i-1\) dimensional kernel.  The M\"obius signs
rotate this decomposition by a diagonal unitary; they do not eliminate
the kernel.  Hence a support-only lower inequality
\begin{equation}
 \|\Eop_\mu z\|^2\ge c\|z\|^2
 \label{eq:tpc58-intro-false-ambient}
\end{equation}
with \(c>0\) is false on the full ambient space whenever any \(k_i\ge2\).
This negative result is exact, but its quantifier is equally important:
it does not show failure on the smaller image traced out by the literal
physical coefficients.

\subsection{The actual coefficient image}

Let a fixed baseline carrier be chosen independently of the bounded row
coefficient vector \(\zeta\), and let
\begin{equation}
 T:\C^{\cM}\longrightarrow\cH_{\rm W},
 \qquad z=T\zeta,
 \qquad E:=\Ran T.
 \label{eq:tpc58-intro-T}
\end{equation}
The best coefficient-uniform comparison is
\begin{equation}
 \alpha_E
 :=\inf_{0\ne z\in E}
   \frac{\|\Eop_\mu z\|^2}{\|z\|^2}.
 \label{eq:tpc58-intro-alpha}
\end{equation}
In finite dimension it is the least eigenvalue, including zero, of the
compression of \(\Eop_\mu^*\Eop_\mu\) to \(E\), and
\begin{equation}
 \alpha_E>0
 \quad\Longleftrightarrow\quad
 E\cap\Ker\Eop_\mu=\{0\}.
 \label{eq:tpc58-intro-transversality}
\end{equation}
With
\begin{equation}
 R:=T^*T,
 \qquad N:=T^*\Eop_\mu^*\Eop_\mu T,
 \qquad \cS:=\Ran R,
 \label{eq:tpc58-intro-two-Grams}
\end{equation}
the exact relative defect is
\begin{equation}
 \Theta:=R^{\dagger/2}(N-R)R^{\dagger/2}\big|_{\cS},
 \qquad
 \alpha_E=1+\lambda_{\min}(\Theta).
 \label{eq:tpc58-intro-relative-defect}
\end{equation}
The restriction to \(\cS\) is essential: zero directions in \(\Ker T\)
must not be counted as physical generalized eigenvectors.

There are therefore three different questions:
\begin{enumerate}[label=(\roman*)]
 \item the ambient question, which has a sharp dark-space obstruction;
 \item the coefficient-uniform question on the actual image \(E\); and
 \item the distinguished-vector question for the one literal coefficient
       vector inherited at scale \(X\).
\end{enumerate}
Failure of (i) does not imply failure of (ii) or (iii).  Conversely, a
numerically favorable distinguished vector does not prove a uniform
bound on \(E\).

\subsection{The fixed-shift two-stage gate}

Use the sign convention
\begin{equation}
 z_H(i,s):=-\sum_{p>y}b_{i,sp},
 \qquad
 h:=\Eop_{\mu,H}z_H=-\sum_{p>y}v_p.
 \label{eq:tpc58-intro-high-convention}
\end{equation}
The TPC--57 terminal residual amplitude already contains the common
physical M\"obius sign \(\mu(s)\).  If that vector is denoted by \(A\),
then the literal interface is
\begin{equation}
 A=\Uop_\mu z_H,
 \qquad
 \Eop_1A=h,
 \qquad
 \|A\|=\|z_H\|.
 \label{eq:tpc58-intro-terminal-interface}
\end{equation}
Thus no second M\"obius factor is inserted when the TPC--57 output is
grouped natively.

For the actual high coefficient image define \(\alpha_H\) by the
analogue of \eqref{eq:tpc58-intro-alpha}.  Separately, on the actual
paired image of \(v_0\) and \(h\), define
\begin{equation}
 \delta_{\rm aff}
 :=\inf_{\substack{(b,g)\in\cF_{\rm aff}\\(b,g)\ne(0,0)}}
   \frac{\|b+g\|^2}{\|b\|^2+\|g\|^2}.
 \label{eq:tpc58-intro-delta}
\end{equation}
Here \(\cF_{\rm aff}\) is the nonzero image of the coupled map driven by
the same row vector in both entries, defined precisely in
\cref{sec:tpc58-affine-completion}.  Then the exact route is
\begin{equation}
 \cN_{\rm nat}
 =\|v_0+h\|^2
 \ge \delta_{\rm aff}\alpha_H\|z_H\|^2.
 \label{eq:tpc58-intro-two-stage}
\end{equation}
The constants are independent gates.  Residual labels may cancel under
native erasure even if the terminal-prime fibers are singletons, and the
base vector may cancel the high vector even if native high energy is
positive.

There is also an upstream support gate.  At the endpoint cutoff
\(y=X^{267/400+\eta}\) with fixed \(\eta>0\), the relation
\(pr=mj+h_0\le X^{1+o(1)}\) forces
\(r\le X^{133/400-\eta+o(1)}\).  Hence this very-high face is eventually
disjoint from the older dyadic block
\(R\asymp J=X^{133/400+o(1)}\).  The interface identities remain exact on
their declared face, but they do not prove that the face carries a positive
fraction of the inherited energy.  A compatible-band relocation and
occupancy estimate is therefore charged separately in the endpoint ledger.

\subsection{Claim boundary and organization}

The finite-dimensional erasure spectrum, compressed Rayleigh quotient,
generalized-Gram formula, and counterexamples are unconditional L0
results.  The coordinate dictionary at one fixed \(h_0\ne0\) is an L1
interface statement on the declared baseline carrier.  A positive
polynomial lower bound for either \(\alpha_H\) or \(\delta_{\rm aff}\),
relocation to and occupancy of the compatible smaller cofactor band,
terminal activation at the endpoint, boundary/tail reconnection, and an
estimate for the original arithmetic defect remain open.  In particular,
no parity improvement and no twin-prime conclusion is asserted; the
classical parity boundary remains in force \citep{FriedlanderIwaniec2010}.

\Cref{sec:tpc58-two-gram-dictionary} fixes the two Gram forms and their
literal signs.  \Cref{sec:tpc58-erasure-spectrum} proves the exact erasure
spectrum, and \cref{sec:tpc58-restricted-range} gives the actual-image
and generalized-Gram theorems.  The fixed-shift coordinate match is
proved in \cref{sec:tpc58-fixed-shift-interface}; the full affine
completion follows in \cref{sec:tpc58-affine-completion}.
\Cref{sec:tpc58-transport-ledger,sec:tpc58-claims-next} record the
transport restrictions, endpoint budget, kill criteria, and next
arithmetic gate.
